\chapter*{\prefacename}
\addcontentsline{toc}{chapter}{\prefacename}

Η παρούσα εργασία γεννήθηκε μέσα από τρία χρόνια συμμετοχής στην ομάδα \en{Energy Transition (ET) Platform} της Maersk, ως μέλος μιας διεθνούς ομάδας μηχανικών λογισμικού που ανέπτυξε ένα σύστημα παραγωγικής κλίμακας για τη μέτρηση και διαχείριση εκπομπών αερίων θερμοκηπίου. Αυτή η εμπειρία ήταν ταυτόχρονα επαγγελματική πρόκληση και πρωτότυπη ευκαιρία: να σχεδιαστεί και να τεθεί σε λειτουργία λογισμικό που συνδέει μηχανική ακρίβεια με έναν στόχο βαθύτατα περιβαλλοντικής σημασίας.

Η επιλογή της Maersk ως μελέτης περίπτωσης δεν ήταν τυχαία. Μια εταιρεία που μεταφέρει περίπου το 17\% των παγκόσμιων εμπορευμάτων κοντέινερ, που δεσμεύτηκε για \en{net zero} εκπομπές το 2040 --- δέκα χρόνια πριν από τον αντίστοιχο στόχο του ΙΜΟ --- και που το 2023 αντιμετώπισε εσωτερικό έλεγχο για την αξιοπιστία των δεδομένων εκπομπών της, αντιπροσωπεύει μια κατηγορία προβλημάτων όπου το λογισμικό δεν είναι απλώς εργαλείο υποστήριξης, αλλά κρίσιμος παράγοντας της επιχειρηματικής και περιβαλλοντικής στρατηγικής.

Η εργασία επιχειρεί να τεκμηριώσει αυτή την εμπειρία με ακρίβεια: όχι απλώς να περιγράψει τι κατασκευάστηκε, αλλά να εξηγήσει γιατί έγιναν συγκεκριμένες αρχιτεκτονικές και μεθοδολογικές επιλογές, ποιες ήταν οι συνέπειές τους, και τι μαθήματα μεταφέρονται πέραν της συγκεκριμένης υλοποίησης.

\begin{flushright}
\textit{Κωνσταντίνος Γκίνης}\\
\textit{Πάτρα, Μάιος 2026}
\end{flushright}
